% Chapter 20: The Sato-Tate Conjecture and the Langlands Program
\let\LocalMacro\relax

\chapter{The Sato-Tate Conjecture and the Langlands Program}

\section{The Problem}
Let $E/\mathbb{Q}$ be an elliptic curve without complex multiplication. For each prime $p$ of good reduction, let $a_p(E) = p + 1 - \#E(\mathbb{F}_p)$. The \emph{Sato-Tate conjecture} (1960) predicts the distribution of the normalized Frobenius traces:

\begin{equation}
\frac{a_p(E)}{2\sqrt{p}} \in [-1, 1]
\end{equation}

\begin{conjecture}[Sato-Tate]
For a non-CM elliptic curve $E/\mathbb{Q}$, the sequence $\left\{ \frac{a_p(E)}{2\sqrt{p}} \right\}_{p \text{ good}}$ is equidistributed in $[-1, 1]$ with respect to the \emph{Sato-Tate measure}
\[
\mu_{\text{ST}} = \frac{2}{\pi} \sqrt{1 - t^2} \, dt.
\]
\end{conjecture}

\begin{historicalbox}
Sato and Tate independently observed this distribution numerically for $y^2 = x^3 - x$ (which has CM) and for non-CM curves like $y^2 = x^3 - 4x^2 + 2x$. The conjecture remained open for 50 years, becoming a cornerstone test for the Langlands program.
\end{historicalbox}

\section{Mathematical Theory}


\subsection{Prerequisites}
This chapter assumes familiarity with:
\begin{itemize}
\item Algebraic number theory (Galois representations, L-functions)
\item Elliptic curves and modular forms
\item Basic representation theory of compact groups
\item The classical Langlands program (reciprocity laws)
\end{itemize}

\subsection{Galois Representations and Automorphic Forms}
For an elliptic curve $E/\mathbb{Q}$, the $p$-adic Tate module $T_p(E)$ gives a Galois representation
\[
\rho_{E,p}: \text{Gal}(\overline{\mathbb{Q}}/\mathbb{Q}) \to \text{GL}_2(\mathbb{Z}_p).
\]
The Sato-Tate conjecture is equivalent to the statement that the restriction of $\rho_{E,p}$ to the decomposition group at $p$ matches a modular form of weight 2.

\subsection{The Role of the Langlands Program}
The Langlands program predicts a correspondence:
\begin{itemize}
\item \textbf{Galois side}: $n$-dimensional Galois representations $\rho: \text{Gal}(\overline{\mathbb{Q}}/\mathbb{Q}) \to \text{GL}_n(\mathbb{C})$.
\item \textbf{Automorphic side}: Automorphic representations $\pi$ of $\text{GL}_n(\mathbb{A}_{\mathbb{Q}})$.
\end{itemize}
For $n=2$, the Sato-Tate conjecture is a special case of the \emph{modularity theorem} (formerly Taniyama-Shimura-Weil) plus the \emph{generalized Sato-Tate conjecture} for higher-genus curves.

\subsection{Equidistribution and Moments}
The Sato-Tate measure $\mu_{\text{ST}}$ has moments:
\[
\int_{-1}^1 t^k \, d\mu_{\text{ST}}(t) = \begin{cases}
0 & k \text{ odd} \\
C_{k/2} & k \text{ even}
\end{cases}
\]
where $C_n = \frac{1}{n+1}\binom{2n}{n}$ are Catalan numbers. Verifying these moments for $a_p(E)$ is equivalent to the conjecture.

\section{The Discovery}

\subsection{Modularity and Potential Modularity}
\begin{itemize}
\item \textbf{Wiles \cite{wiles1995}}: Modularity for semistable elliptic curves (proved FLT).
\item \textbf{Breuil-Conrad-Diamond-Taylor \cite{bcdt2001}}: Full modularity for all elliptic curves over $\mathbb{Q}$.
\item \textbf{Taylor (2002)}: Potential modularity for elliptic curves over totally real fields.
\end{itemize}

\subsection{The Proof (2008-2011)}
The Sato-Tate conjecture was proved in a series of papers by \textbf{Clozel, Harris, Shepherd-Barron, and Taylor} (2008, 2011).

\begin{theorem}[Sato-Tate \cite{barnetlamb2011}]
For any non-CM elliptic curve $E/\mathbb{Q}$, the Sato-Tate conjecture holds.
\end{theorem}

The proof strategy:
\begin{enumerate}
\item \textbf{Potential automorphy}: Show that the Galois representation $\rho_{E,\lambda}$ is potentially automorphic (becomes automorphic after a solvable base change).
\item \textbf{Stable trace formula}: Use Arthur's stable trace formula for $\text{GL}_2$ to relate Galois representations to automorphic forms on $\text{GL}_2$ over CM fields.
\item \textbf{Equidistribution}: Transfer the equidistribution of Frobenius traces from the automorphic side (where it follows from the spectral theorem) to the Galois side.
\item \textbf{CM case}: The CM case was already known by Deuring (1941) and uses Hecke characters.
\end{enumerate}

\begin{highlightbox}
The proof is a landmark of the Langlands program: it uses the full power of the stable trace formula, base change, and potential automorphy to transfer a Galois-theoretic question to an automorphic one where harmonic analysis applies.
\end{highlightbox}

\subsection{Generalized Sato-Tate}
The conjecture generalizes to higher-genus curves and higher-dimensional abelian varieties. The equidistribution measure is determined by the \emph{Sato-Tate group}, a compact subgroup of $\text{USp}(2g)$.

\section{Impact}
\begin{itemize}
\item \textbf{Langlands program}: First major application of the stable trace formula to a classical number-theoretic problem.
\item \textbf{Modularity}: Extended modularity techniques to non-semistable curves and higher-dimensional varieties.
\item \textbf{Computational number theory}: Verified Sato-Tate distributions for millions of curves.
\item \textbf{Arithmetic statistics}: Led to precise conjectures on the distribution of $a_p$ for higher-genus curves (FKRS conjecture).
\end{itemize}

\section{Exercises}

\begin{exercise}[Basic]
For $E: y^2 = x^3 - x$ (CM by $\mathbb{Z}[i]$), compute the first 10 values of $a_p$ and explain why they do \emph{not} follow the Sato-Tate distribution.
\end{exercise}

\begin{exercise}[Intermediate]
Show that the Sato-Tate measure $\frac{2}{\pi}\sqrt{1-t^2} dt$ is the eigenvalue distribution of a random matrix in $\text{SU}(2)$ with Haar measure.
\end{exercise}

\begin{exercise}[Challenge]
Explain why the stable trace formula is needed instead of the ordinary trace formula to prove Sato-Tate.
\end{exercise}

\section{Timeline}

\begin{tabularx}{\linewidth}{lX}
\toprule
\textbf{Year} & \textbf{Event} \\ \midrule
1960 & Sato and Tate independently conjecture the distribution \\ 
1968 & Serre proves Sato-Tate for CM elliptic curves \\ 
1995 & Wiles proves modularity for semistable curves \\ 
2001 & BCDT prove modularity for all elliptic curves over $\mathbb{Q}$ \\ 
2008 & Clozel-Harris-Shepherd-Barron-Taylor announce Sato-Tate proof \\ 
2011 & Full proof published in JAMS \\ 
2010s & Generalized Sato-Tate for higher-genus curves (FKRS) \\ \bottomrule
\end{tabularx}

\section{Citations}

\begin{enumerate}
\item Clozel, L., Harris, M., \& Taylor, R. (2008). ``Automorphy for some $\ell$-adic lifts of automorphic mod $\ell$ Galois representations.'' Publications Mathématiques de l'IHÉS, 108(1), 1--181.
\item Barnet-Lamb, T., Geraghty, D., Harris, M., \& Taylor, R. (2011). ``A family of Calabi-Yau varieties and potential automorphy II.'' Publications Mathématiques de l'IHÉS, 116(1), 119--180.
\item Serre, J.-P. (1968). ``Abelian $\ell$-adic representations and elliptic curves.'' W. A. Benjamin.
\item Harris, M. (2011). ``The Sato-Tate conjecture.'' Notices of the AMS, 58(8), 1046--1056.
\end{enumerate}
